\subsection{同伦的连续映射}

定义 1 —— 设 X 与 Y 是拓扑空间，f 与 g 是从 X 到 Y 的连续映射。连接 f 与 g 的同伦是一个连续映射 \(\sigma: X \times I \to Y\)，使得对每个 \(x \in X\)，有 \(\sigma(x,0) = f(x)\) 与 \(\sigma(x,1) = g(x)\)。若存在连接 f 与 g 的同伦，则称 f 与 g 同伦。

我们称 \(f\) 为起点，\(g\) 为同伦 \(\sigma\) 的终点（参见下文 III, p. 257, 注 2）。

设 A 是 X 的一个子集。若同伦 \(\sigma\) 满足：对每个 \(a \in \mathrm{A}\)，从 I 到 Y 的映射 \(t \mapsto \sigma(a, t)\) 是常值的，则称它在 A 上是固定的。此时，\(\sigma\) 的起点与终点在 A 的每点处取相同的值。

设 X 与 Y 是拓扑空间。我们称同伦 \(\sigma\colon X\times I\to Y\) 与 \(\tau\colon X\times I\to Y\) 是可拼接的，如果 \(\sigma\) 的终点是 \(\tau\) 的起点，亦即有 \(\sigma(x,1)=\tau(x,0)\)（对所有 \(x\in X\)）。此时，映射 \(\sigma*\tau\)（从 \(X\times I\) 到 Y）由下式定义：

\[
(\sigma * \tau) (x, t) = \left\{ \begin{array}{l l} \sigma (x, 2 t) & \text {for} 0 \leqslant t \leqslant 1 / 2 \\ \tau (x, 2 t - 1) & \text {for} 1 / 2 \leqslant t \leqslant 1 \end{array} \right.\tag{1}
\]

它是连续的（TG, I, p. 19, 命题 4），且是连接 \(\sigma\) 的起点与 \(\tau\) 的终点的一个同伦。称之为同伦 \(\sigma\) 与 \(\tau\) 的拼接同伦。

若 \(\sigma\colon X\times\mathbf{I}\to Y\) 是一个同伦，则映射 \(\overline{\sigma}\colon X\times\mathbf{I}\to Y\)（由 \((x,t)\mapsto\sigma(x,1-t)\) 定义）是连接 \(\sigma\) 的终点与 \(\sigma\) 的起点的同伦。我们有 \(\overline{\overline{\sigma}}=\sigma\)。若 \(\sigma\) 与 \(\tau\) 是从 \(X\times\mathbf{I}\) 到 Y 的可拼接同伦，则同伦 \(\overline{\tau}\) 与 \(\overline{\sigma}\) 可拼接，且有 \(\overline{\sigma*\overline{\tau}}=\overline{\tau}*\overline{\sigma}\)。

命题 1 —— 设 X 与 Y 是拓扑空间。关系“f 与 g 同伦”是从 X 到 Y 的连续映射所成的集合 \(\mathcal{C}(\mathrm{X};\mathrm{Y})\) 上的一个等价关系。

设 \(f\) 是 \(\mathcal{C}(\mathrm{X};\mathrm{Y})\) 的一个元素。映射 \(f\circ\mathrm{pr}_{1}\colon\mathrm{X}\times\mathrm{I}\to\mathrm{Y}\) 是连接 \(f\) 与 \(f\) 的同伦；因此该关系是自反的。

设 \(f\) 与 \(g\) 是 \(\mathcal{C}(\mathrm{X};\mathrm{Y})\) 的元素，\(\sigma\colon\mathrm{X}\times\mathbf{I}\to\mathrm{Y}\) 是连接 \(f\) 与 \(g\) 的同伦。于是映射 \(\overline{\sigma}\colon\mathrm{X}\times\mathbf{I}\to\mathrm{Y}\) 是连接 \(g\) 与 \(f\) 的同伦；因此所考虑的关系是对称的。

最后证明它是传递的。若 \(f\)、\(g\)、\(h\) 是 \(\mathcal{C}(\mathrm{X};\mathrm{Y})\) 的元素，\(\sigma\) 是连接 \(f\) 与 \(g\) 的同伦，\(\tau\) 是连接 \(g\) 与 \(h\) 的同伦，则 \(\sigma\) 与 \(\tau\) 可拼接，且 \(\sigma*\tau\) 是连接 \(f\) 与 \(h\) 的同伦。

设 X 与 Y 是拓扑空间。\(\mathcal{C}(\mathrm{X};\mathrm{Y})\) 中的等价关系“f 与 g 同伦”（命题 1）称为同伦关系。\(\mathcal{C}(\mathrm{X};\mathrm{Y})\) 对该关系的商集记作 {[}X;Y{]}。它的元素称为从 X 到 Y 的连续映射的同伦类。连续映射 \(f\colon \mathrm{X}\to \mathrm{Y}\) 的同伦类常记作 {[}f{]}。

命题 2 —— 设 X、Y、Z 是拓扑空间，f 与 \(f'\) 是从 X 到 Y 的连续映射，g 与 \(g'\) 是从 Y 到 Z 的连续映射。若 f 与 \(f'\) 同伦且 g 与 \(g'\) 同伦，则 \(g \circ f\) 与 \(g' \circ f'\) 同伦。

设 \(\sigma\) 是连接 \(f\) 与 \(f'\) 的同伦，\(\tau\) 是连接 \(g\) 与 \(g'\) 的同伦。于是映射 \(\theta\colon X\times I\to Z\)（由 \(\theta(x,t)=\tau(\sigma(x,t),t)\) 定义）是连接 \(g\circ f\) 与 \(g'\circ f'\) 的同伦。

设 X、Y、Z 是拓扑空间。给定同伦类 \(\varphi \in [\mathrm{X};\mathrm{Y}]\) 与 \(\psi \in [\mathrm{Y};\mathrm{Z}]\)，诸映射 \(g \circ f: \mathrm{X} \to \mathrm{Z}\)（其中 \(f \in \varphi\)，\(g \in \psi\)）都属于同一个同伦类（命题 2），把它记作 \(\psi \circ \varphi\)，并称为类 \(\psi\) 与 \(\varphi\) 的复合同伦类。映射 {[}X;Y{]} \(\times\) {[}Y;Z{]} \(\to\) {[}X;Z{]}（它把 \(\psi \circ \varphi\) 赋予 \((\varphi,\psi)\)）称为合成映射。

设 \(\varphi\in[\mathrm{X};\mathrm{Y}]\)；我们有 \(\varphi=\varphi\circ[\mathrm{Id}_{\mathrm{X}}]=[\mathrm{Id}_{\mathrm{Y}}]\circ\varphi\)。

设 X、Y、Z、T 是拓扑空间，\(\varphi \in [\mathrm{X};\mathrm{Y}]\)、\(\psi \in [\mathrm{Y};\mathrm{Z}]\)、\(\chi \in [\mathrm{Z};\mathrm{T}]\)。同伦类 \(\chi \circ (\psi \circ \varphi)\) 等于 \((\chi \circ \psi) \circ \varphi\)；把它记作 \(\chi \circ \psi \circ \varphi\)。

命题 3 —— 设 X 是拓扑空间，\((\mathrm{Y}_{j})_{j\in\mathrm{J}}\) 是一族拓扑空间。从 \([X;\prod_{j\in\mathrm{J}}Y_{j}]\) 到积集 \(\prod_{j\in\mathrm{J}}[X;Y_{j}]\) 的、由 \(\varphi\mapsto([\mathrm{pr}_{j}]\circ\varphi)_{j\in\mathrm{J}}\) 定义的映射是双射。

满射性由 I, p. 25, 命题 1 立得。我们来证明单射性。设 \(f\) 与 \(g\) 是从 X 到 \(\prod_{j \in \mathrm{J}} \mathrm{Y}_j\) 的连续映射。对所有 \(j \in \mathrm{J}\)，令 \(f_j = \mathrm{pr}_j \circ f\)，\(g_j = \mathrm{pr}_j \circ g\)；设 \(f_j\) 与 \(g_j\) 同伦，\(\sigma_j\) 是连接 \(f_j\) 与 \(g_j\) 的同伦。映射 \(\sigma = (\sigma_j)\)（从 \(\mathrm{X} \times \mathbf{I}\) 到 \(\prod_{j \in \mathrm{J}} \mathrm{Y}_j\)）是连续的（出处同上）；它是连接 \(f\) 与 \(g\) 的同伦，由此得命题。

推论 —— 设 \((\mathrm{X}_{j})_{j\in\mathrm{J}}\) 与 \((\mathrm{Y}_{j})_{j\in\mathrm{J}}\) 是有同一指标集的拓扑空间族。对每个 \(j\in J\)，设 \(f_{j}\) 与 \(g_{j}\) 是从 \(X_{j}\) 到 \(Y_{j}\) 的连续映射。若对每个 \(j\in J\)，映射 \(f_{j}\) 与 \(g_{j}\) 都同伦，则积映射 \(f\colon(x_{j})\mapsto(f_{j}(x_{j}))\) 与 \(g\colon(x_{j})\mapsto(g_{j}(x_{j}))\)（从 \(\prod_{j\in J}X_{j}\) 到 \(\prod_{j\in J}Y_{j}\)）也同伦。

假设有 \([f_j] = [g_j]\)（对所有 \(j \in \mathbf{J}\)）。我们有 \([\mathrm{pr}_j] \circ [f] = [f_j \circ \mathrm{pr}_j]\) 与 \([\mathrm{pr}_j] \circ [g] = [g_j \circ \mathrm{pr}_j]\)；于是由命题 2 得 \([\mathrm{pr}_j] \circ [f] = [\mathrm{pr}_j] \circ [g]\)，再由命题 3 得 \([f] = [g]\)。

